\chapter{Face Masks}
\index{Face Masks}

\section*{Definition and Overview}
Face masks are coverings worn over the nose and mouth that reduce the spread of respiratory infections. They work in two ways: they catch the droplets and particles exhaled by the wearer, protecting others, and they filter the air the wearer breathes, protecting the wearer from infection. Face masks were widely used during the COVID-19 pandemic and remain important for people with respiratory infections, in healthcare settings, and for those at higher risk. There are different types of masks with different levels of protection, and correct wearing is essential for effectiveness.

\section*{Epidemiology}
Respiratory infections spread mainly through droplets and airborne particles released when people cough, sneeze, talk, and breathe. During the COVID-19 pandemic, mask wearing was shown to reduce transmission in the community and in healthcare settings. Masks continue to be recommended for people with symptoms, in high-risk settings, and for people vulnerable to severe illness, particularly during outbreaks of influenza, RSV, and COVID-19.

\section*{Aetiology and Pathophysiology}
Masks work by filtering and containing respiratory particles.
\begin{itemize}
    \item \textbf{Droplets and aerosols:} respiratory infections spread through large droplets and fine aerosols
    \item \textbf{Source control:} a mask worn by an infected person contains their droplets, protecting others
    \item \textbf{Protection:} a well-fitted mask worn by a healthy person filters particles from the air
    \item \textbf{Types:} cloth masks, surgical masks, and respirators such as N95/P2 masks, with increasing levels of filtration
    \item \textbf{Fit:} a mask only works if it fits snugly over the nose and mouth
    \item \textbf{Hand hygiene:} masks work best with hand washing and other measures
\end{itemize}

\section*{Clinical Features}
Different masks suit different situations.
\begin{itemize}
    \item \textbf{Surgical masks:} loose-fitting, protect others from the wearer's droplets, and give some protection to the wearer
    \item \textbf{Cloth masks:} provide basic source control, with variable filtration
    \item \textbf{N95/P2 respirators:} close-fitting masks that filter fine particles and give the wearer the best protection
    \item \textbf{When to wear:} with respiratory symptoms, in healthcare and crowded settings, and for people at high risk
    \item \textbf{Fitting:} the mask should cover the nose and mouth with no gaps
\end{itemize}

\section*{Differential Diagnosis}
Masks are one measure among several.
\begin{itemize}
    \item Vaccination provides the strongest protection against severe illness
    \item Physical distancing, ventilation, and hand hygiene reduce transmission
    \item Staying home when unwell is the most important step
    \item The measures work best in combination
\end{itemize}

\section*{When to Seek Medical Care}
People with respiratory symptoms who are unwell, at high risk, or developing breathing difficulty should seek medical advice. Masks do not replace medical care for illness.

\textbf{Contact your local emergency services for:}
\begin{itemize}
    \item Severe difficulty breathing
    \item Chest pain
    \item Collapse or confusion
    \item Bluish lips or face
\end{itemize}

\section*{Diagnosis and Investigations}
Choosing and using a mask is practical.
\begin{itemize}
    \item \textbf{Choose the right type:} respirators for high-risk settings, surgical masks for general use
    \item \textbf{Check the fit:} the mask covers the nose, mouth, and chin with no gaps
    \item \textbf{Replace regularly:} masks should be changed when damp, damaged, or as advised
    \item \textbf{Hand hygiene:} wash hands before putting on and after removing a mask
    \item \textbf{Correct removal:} remove by the ear loops without touching the front
\end{itemize}

\section*{Management}
Masks are managed as part of infection control.
\begin{itemize}
    \item \textbf{People with symptoms:} wear a mask when leaving home and when around others
    \item \textbf{High-risk settings:} wear a mask in hospitals, aged care, and crowded indoor spaces
    \item \textbf{Vulnerable people:} a well-fitted respirator offers the best protection
    \item \textbf{Healthcare workers:} follow workplace requirements
    \item \textbf{Outbreaks:} follow public health advice
\end{itemize}

\section*{Complications}
\begin{itemize}
    \item Discomfort, warmth, and difficulty wearing masks for long periods
    \item Skin irritation on the face
    \item A false sense of security if masks are used without other measures
    \item Improper use, including uncovered noses, which removes the protection
\end{itemize}

\section*{Prognosis and Outcomes}
Masks are an effective, simple measure that reduces the spread of respiratory infections when used correctly. Combined with vaccination, ventilation, and hygiene, they protect both individuals and the community. Wearing a mask when unwell is a considerate act that protects others, particularly the vulnerable.

\section*{Prevention and Self-Care}
\begin{itemize}
    \item Wear a mask when you have respiratory symptoms and need to leave home
    \item Use a well-fitted mask in healthcare and crowded settings when advised
    \item Choose a respirator if you are at higher risk of severe illness
    \item Replace masks regularly and dispose of them safely
    \item Combine masks with vaccination, hand washing, and staying home when unwell
    \item Follow public health advice during outbreaks
\end{itemize}

\section*{Sources and Further Reading}
The following reputable sources provide further information for healthcare students and patients seeking reliable, current advice about face masks.
\begin{itemize}
    \item Department of Health and Aged Care. \textit{Masks and respiratory infections.} https://www.health.gov.au/
    \item Healthdirect Australia. \textit{How to use a face mask.} https://www.healthdirect.gov.au/
    \item NSW Health. \textit{Face mask advice.} https://www.health.nsw.gov.au/
    \item World Health Organization. \textit{Masks and respiratory hygiene.} https://www.who.int/
    \item Centers for Disease Control and Prevention. \textit{Mask use and protection.} https://www.cdc.gov/
    \item Therapeutic Goods Administration. \textit{Surgical masks and respirators.} https://www.tga.gov.au/
\end{itemize}
